\chapter{绪\hspace{6pt}论}

\section{研究背景与意义}

Shell（如Bash、Zsh）是Linux/Unix操作系统的核心交互工具，
用户通过命令行直接控制文件系统、进程管理、网络通信等底层功能。
然而，Shell命令语法复杂，涉及路径规则、参数组合、管道与重定向等专业知识，
对非技术背景用户存在较高的使用门槛，
高危命令（如\texttt{rm -rf /}、\texttt{dd of=/dev/sda}）的误执行可造成不可逆的数据损失。

近年来，以GPT系列、LLaMA、Qwen\citing{bai2023qwen}为代表的大语言模型在自然语言理解与代码生成领域取得了显著突破。
将LLM能力与Shell操作结合，构建自然语言驱动的Shell命令生成系统，
是降低操作系统使用门槛、提升运维效率的重要研究方向。
用户只需以中文口语化描述操作意图（如“查看当前目录下超过100MB的文件”），
系统即可自动生成对应的Shell命令（如\texttt{find ./ -size +100M}）并辅助执行。

然而，通用LLM直接应用于Shell场景面临三项核心挑战：
其一，LLM训练数据以英文为主，对中文口语化表达的理解存在偏差；
其二，LLM训练语料中Shell领域长尾知识（特定发行版命令、企业内部操作流程）
分布稀疏，导致模型易产生语法正确但语义错误的命令（即“幻觉”问题）；
其三，Shell命令执行具有强副作用，安全性保障是部署LLM辅助系统的前提条件。

基于上述背景，本课题围绕中文Shell智能体的可靠性与安全性展开研究，
具有重要的理论意义与实践价值。

\section{国内外研究现状}

\subsection{自然语言到代码生成与LLM Agent}

自然语言到代码生成（Natural Language to Code，NL2Code）是人工智能领域的经典研究课题。
早期工作基于统计机器翻译方法将自然语言映射到程序语言，
近年来基于Transformer架构的预训练语言模型大幅提升了代码生成的准确率。
Codex\citing{chen2021evaluating}、Code LLaMA\citing{roziere2023code}等专用代码生成模型的发布标志着该领域的重要进展。
Shell命令生成作为NL2Code的子方向\citing{lin2018nl2bash}，
具有命令结构简短、参数组合多样、中文口语表达形式不固定等特点，
对模型的语义理解能力提出了更高要求。

同时，LLM在推理与工具使用中的能力已成为新的研究焦点。
ReAct框架\citing{yao2022react}通过交织推理步骤（Thought）与行动步骤（Action）的方式，
显著提升了LLM在复杂任务中的成功率。
Toolformer\citing{schick2023toolformer}证实了大语言模型能够自学使用API和工具。
OS-Copilot\citing{lu2023oscopilot}针对操作系统级任务构建了通用的计算机Agent框架，
为本文的系统设计提供了重要参考。
AgentBench\citing{zeng2024agentbench}等评估框架的推出为Agent系统的性能衡量提供了标准化指标。

\subsection{检索增强生成技术}

检索增强生成（Retrieval-Augmented Generation，RAG）通过在推理时动态检索外部知识库
并将相关文档注入提示上下文，有效缓解了LLM的知识局限与幻觉问题。
RAG技术已在开放域问答、文档摘要、代码补全等多个场景中得到广泛验证\citing{lewis2020retrieval,gao2023ragsurvey}。
Self-RAG\citing{asai2023selfrag}进一步引入了检索反思机制，使模型能够动态决定何时检索。
向量检索（基于稠密嵌入的余弦相似度匹配）与稀疏检索（基于BM25等词频统计）
的混合策略被证明在多个评测基准上优于单一检索方式\citing{robertson2009bm25}，
能够在语义相似度与关键词精确匹配之间实现互补。

LLM生成内容的安全性已成为系统部署的重要考量。
在Shell命令执行场景下，提示注入攻击（Prompt Injection）\citing{greshake2023injection}、
危险命令生成等风险需通过系统级防护机制加以规避。
现有工作主要从输出过滤、提示设计和权限控制三个维度进行防护，
但针对Shell命令预执行模拟校验的系统性方案仍相对欠缺。

针对LLM对齐与安全，通过人类反馈的强化学习（RLHF）\citing{ouyang2022training}已成为主流方案。
深度强化学习从人类偏好\citing{christiano2017deep}出发，指导模型学习人类价值观。
最近的工作从提示设计与输出过滤角度探索了针对LLM集成应用的安全防护方案。

工具学习是扩展LLM能力的重要方向。
Tool Augmented Language Models（TALM）\citing{parisi2022talm}使模型能够学习何时调用工具。
API-Bank\citing{ma2023api}基准为工具调用能力提供了系统化评估。

链式思维提示（Chain-of-Thought）\citing{wei2022chain}通过多步推理显著提升了LLM的性能。
大语言模型作为零样本推理器\citing{kojima2022large}展示了其强大的泛化能力。
涌现能力分析\citing{wei2022emergent}揭示了LLM在模型规模增大时的非线性性能跳跃
在Shell命令执行场景下，提示注入攻击（Prompt Injection）\citing{greshake2023injection}、
危险命令生成等风险需通过系统级防护机制加以规避。
现有工作主要从输出过滤、提示设计和权限控制三个维度进行防护，
但针对Shell命令预执行模拟校验的系统性方案仍相对欠缺。

综上，现有工作尚未针对中文Shell场景提供一个同时解决
领域知识覆盖不足、中文语义理解偏差与命令执行安全三项挑战的端到端系统。
本文的研究正是针对这一空白展开。

\section{研究内容与贡献}

本文的主要研究内容与贡献如下：

\begin{enumerate}
    \item 设计并实现了基于LLM的中文Shell智能体系统。
    采用客户端-服务端分离架构，以本地量化Qwen-7B模型为推理核心，
    实现自然语言描述到Shell命令的端到端生成，
    并适配Bash与Zsh两种主流Shell的语法差异。
    系统支持多会话管理机制，允许用户在同一进程中管理多个独立的任务上下文，
    每个会话拥有唯一的标识符与隔离的对话记忆，
    实现了用户在不同工作流之间快速切换的能力。

    \item 构建了面向Shell领域的RAG检索模块。
    以ChromaDB向量数据库与BAAI/bge-small-zh-v1.5中文嵌入模型为基础，
    提出混合检索与关键词路由策略，
    在80个自然语言查询评估集上将Recall@1提升至0.90、MRR提升至0.9375。

    \item 设计并实现了三级命令安全执行框架。
    包含基于白名单/黑名单的权限分级机制、
    预执行“语法校验—风险评估—效果模拟”三步流水线
    及结构化执行日志系统，
    通过58个单元测试与集成测试用例验证其正确性。

    \item 构建了分层对话记忆系统。
    基于SQLite实现短期记忆、长期记忆与摘要记忆三层架构，
    支持多轮对话的上下文关联，
    使系统能够理解如“将这些文件移动到/data/backup”此类依赖历史上下文的后续指令。

    \item 实现了面向复杂任务的分步执行机制。
    基于关键词启发式规则自动检测多步骤任务请求，
    通过逐步生成—安全校验—用户确认—工作记忆更新的循环，
    驱动复杂任务的连续执行，支持最多12步操作序列，
    每步均享有独立的安全保障与可中断能力。
\end{enumerate}

\section{论文组织结构}

本文其余部分组织如下：
第二章介绍与本研究相关的背景知识，
包括检索增强生成技术、向量检索与嵌入模型、混合检索与重排序方法及Shell命令安全机制；
第三章描述系统整体架构及各核心模块的设计，
重点阐述CLI客户端的多会话管理机制、响应解析容错策略、分步任务执行流程及模型服务器的会话记忆管理；
第四章详述领域RAG检索系统的知识库构建、嵌入索引、混合检索策略与关键词路由过滤；
第五章介绍命令安全执行系统的权限分级模型、预执行校验流水线及执行日志系统；
第六章展示实验设置、评估指标与消融实验结果，并进行分析与讨论；
第七章总结全文主要工作与创新点，并对未来工作进行展望。

\section{本章小结}

本章阐述了以自然语言驱动Shell命令生成的研究背景与意义，
指出通用LLM在中文Shell场景下面临的中文语义理解偏差、领域知识覆盖不足与命令执行安全三项核心挑战，
回顾了自然语言到代码生成、检索增强生成及LLM安全领域的相关研究现状，
并明确了本文的研究内容与主要贡献，包括多会话隔离机制、混合检索策略、分步执行框架等创新点。
后续各章将围绕系统架构设计、多会话管理、分层记忆系统、RAG检索策略、命令安全执行与实验验证逐步展开。



